\documentclass[a4paper, 12pt]{article}
\usepackage[T2A]{fontenc}
%\usepackage[russian, english]{babel} FUCK IT!
\usepackage[english,russian]{babel}
\usepackage[utf8]{inputenc}
\usepackage{csquotes}
\usepackage{amsmath}
\usepackage{amsfonts}
\usepackage{amssymb}
\usepackage{graphicx}
\usepackage{tikz, transparent, fancybox}
\usetikzlibrary{trees}
\usetikzlibrary{positioning,decorations.pathreplacing,shapes}
\usepackage{smartdiagram}
\usetikzlibrary{arrows,automata}
\usetikzlibrary{positioning}
\usetikzlibrary{shapes,arrows}
\usetikzlibrary{calc, shapes, shadows, arrows, patterns}
\tikzstyle{line} = [ draw, -latex']  

\def\RR{\mathbb{R}}  
\newcommand{\geqs}{\geqslant}
\newcommand{\leqs}{\leqslant}
\newcommand{\X}{{\mathcal X}}
\newcommand{\R}{{\mathcal R}}
\newcommand{\Iref}{I_\text{ref}}
\newcommand{\Iinc}{I_\text{inc}}
\newcommand{\vphi}{\varphi}
\newcommand{\ol}{\overline}
\newcommand{\norm}[1]{\left\|#1\right\|}
\newcommand\myeq{\mathrel{\stackrel{\makebox[0pt]{\mbox{\normalfont\tiny def}}}{=}}}



\textwidth 165mm 
\textheight 240mm
\renewcommand\baselinestretch{1.25}
\topmargin -16mm
\oddsidemargin 0pt
\evensidemargin 0pt

%\usepackage{indentfirst}
%\frenchspacing

            
            
\begin{document}

%\begin{titlepage}
%\newcommand{\HRule}{\rule{\linewidth}{0.5mm}} % Defines a new command for the horizontal lines, change thickness here
%\center 
%
%%----------------------------------------------------------------------------------------
%%   HEADING SECTIONS
%%----------------------------------------------------------------------------------------
%\textsc{Московский Государственный Университет \\ им. М.В.Ломоносова}\\[1.5cm] % Name of your university/college
%
%\includegraphics[width=0.2\textwidth]{pictures/logo.png} \\
%\vfill
%\textsc{Физический факультет}\\[0.5cm]
%\textsc{Кафедра математического моделирования и информатики}\\[0.5cm]
%%----------------------------------------------------------------------------------------
%%   TITLE SECTION
%%----------------------------------------------------------------------------------------
%\vfill
%{ \Huge\bfseries Нейросетевая реализация морфологического метода идентификации объектов по изображениям	}\\[0.4cm] % Title of your document
%\vfill
%%----------------------------------------------------------------------------------------
%%   AUTHOR SECTION
%%----------------------------------------------------------------------------------------
%\begin{minipage}{0.3\textwidth}
%\begin{flushleft} \large
%\emph{Выполнил:}\\
%Байков В.Г.
%\end{flushleft}
%\end{minipage}
%~
%\begin{minipage}{0.65\textwidth}
%\begin{flushright} \large
%\emph{Научный руководитель:} \\
%д.ф.-м.н. проф. Чуличков А.И.\\
%\emph{Соруководитель:} \\
%м.н.с. Зубюк А.В.
%\end{flushright}
%\end{minipage}\\[3cm]
%
%% If you don't want a supervisor, uncomment the two lines below and remove the section above
%%\Large \emph{Author:}\\
%%John \textsc{Smith}\\[3cm] % Your name
%
%%----------------------------------------------------------------------------------------
%%   DATE SECTION
%%----------------------------------------------------------------------------------------
%\vfill % Fill the rest of the page with whitespace
%
%{\large Москва 2017}\\ % Date, change the \today to a set date if you want to be precise
%\end{titlepage}

\begin{titlepage}

\begin{center}
\vfill

{\small ФЕДЕРАЛЬНОЕ ГОСУДАРСТВЕННОЕ БЮДЖЕТНОЕ ОБРАЗОВАТЕЛЬНОЕ 
УЧРЕЖДЕНИЕ ВЫСШЕГО ОБРАЗОВАНИЯ 
«МОСКОВСКИЙ ГОСУДАРСТВЕННЫЙ УНИВЕРСИТЕТ» имени М.В. ЛОМОНОСОВА
\ \\[1ex]
ФИЗИЧЕСКИЙ ФАКУЛЬТЕТ МГУ }
\ \\[0.8ex]
Кафедра математического моделирования и информатики
\ \\[1.8ex] 
\hrule 

\vfill

{\large\bf Нейросетевая реализация морфологического метода идентификации объектов по изображениям \\}
\ \\[2ex]
\begin{flushright}
\begin{minipage}{0.35\textwidth}
{\bf Дипломная работа}  \\ студента 435 группы \\ Байкова~В.~Г.
\\[1ex]
{\bf Научный руководитель} \\ Чуличков~А.~И.\\
{\bf Соруководитель} \\ м.н.с Зубюк~А.~В.
\end{minipage}
\end{flushright}

\vfill
\end{center}
\hspace{3em}\begin{minipage}{0.65\textwidth}
\end{minipage}

\vfill
\begin{center}
Москва, 2017
\end{center}

\end{titlepage}


\newpage
\pagestyle{empty}
\tableofcontents

\newpage
\pagestyle{plain}
\section{Введение}
Трудности при решении многих задач анализа изображений и сигналов связаны с тем, что зарегистрированные изображения и сигналы помимо полезной для исследователя информации несут также информацию об условиях их регистрации, которые исследователю не известны и могут неконтролируемо изменяться. Так, например, изображение какого-либо объекта несёт информацию не только о его геометрической форме, ориентации в пространстве и других характеристиках, присущих непосредственно этому объекту, но также и об условиях освещения, при которых получено данное изображение, характеристиках и настройках использованной камеры и т.\,п. Пусть решается задача узнавания изображённого объекта, тогда условия освещения и характеристики камеры следует считать <<мешающими>> параметрами, изменение которых не должно влиять на результат узнавания. Приведём ещё один пример. Порождённый удалённым источником~--- таким как ядерный взрыв или горный обвал~--- акустический сигнал, распространяясь в атмосфере, претерпевает значительные искажения, вследствие чего несёт информацию не только об источнике звука, но и об акустических свойствах атмосферы. При решении задачи идентификации источника звука <<мешающими>> параметрами следует считать акустические свойства атмосферы, и их изменение не должно влиять на результат идентификации.

Для решения подобных задач были разработаны методы морфологического анализа изображений и сигналов, в рамках которых изображения и сигналы различной природы рассматриваются как элементы функциональных пространств, что позволяет использовать для их описания и анализа универсальные методы линейной алгебры и функционального анализа. В связи с этим далее в настоящей работе будем говорить об анализе изображений, подразумевая, что те же методы могут быть применены и для анализа различных сигналов. В основе методов морфологического анализа лежит идея инвариантности относительно изменения неизвестных и неконтролируемых условий регистрации. Максимальный инвариант класса всех таких изменений называется \emph{формой изображения}\cite{Pitiev_Zadachi_Morph_Analiza}, она может быть определена как множество всех изображений, несущих одну и ту же полезную информацию и отличающихся лишь условиями их регистрации. При решении прикладных задач методами морфологического анализа используется техника проецирования на форму изображения.

Другим распространённым подходом к решению некоторых задач интерпретации изображений и сигналов, таких как идентификация объектов по их изображениям или идентификация типов источников звука по зарегистрированным акустическим сигналам, является применение различных методов машинного обучения

В настоящей работе рассмотрена модель регистрации изображений объектов, освещаемых некогерентными источниками света; при этом предполагается, что каждый светочувствительный детектор используемой для регистрации камеры даёт сигнал на выходе, пропорциональный падающему на него световому потоку. В этом случае множество всех изображений объекта представляет собой выпуклый замкнутый конус. В дипломной работе сконструирована и исследована нейронная сеть, осуществляющая проецирование на форму изображения, заданную таким конусом. В работе приведены результаты тестирования разработанной сети на реальных изображениях.

\section{Рассматриваемая модель}
\subsection{Модель регистрации изображений}
\mbox{}\vspace{-\baselineskip}

\emph{Полем зрения} будем называть конечное подмножество $\X=\{x_1,\, \ldots,\, x_n\}$ координатной плоскости $(-\infty,\infty)\times(-\infty,\infty)$, состоящее из точек $x_1,\, \ldots,\, x_n$, называемых пикселями. \emph{Изображением} будем называть всякую функцию $f: \X\to(-\infty,\infty)$, значения $f(x)$ которой будем называть \emph{яркостями изображения} $f$, $x\in\X$. Линейную комбинацию изображений $f$ и $g$ с коэффициентами $\alpha$ и $\beta$ соответственно определим как изображение
$$
    x\mapsto \alpha f(x) + \beta g(x),
$$
а их скалярное произведение $(f,g)$~--- как число
$$
    (f,g) = \sum_{x\in\X} f(x)g(x).
$$
Множество всех изображений с определёнными выше линейными операциями и скалярным произведением образует евклидово пространство размерности $n$, которое будем называть \emph{пространством изображений} и обозначать $\R$. \cite{PitievChulichkov_morph_methods}, \cite{Pitiev_Zadachi_Morph_Analiza}

\subsection{Форма изображения объекта, регистрируемого при неизменном ракурсе}
\mbox{}\vspace{-\baselineskip}

Рассмотрим регистрацию изображений одного и того же объета при неизменном ракурсе съёмки \cite{ZubukFedotov}. Совокупность скалярных параметров, определяющих выбранный ракурс, обозначим $\vphi$. Такими параметрами могут быть, например, два угла, задающих ориентацию снимаемого объекта в пространстве. Варьироваться могут лишь условия освещения объекта. Сформулированные выше ограничения позволяют определить, что представляет собой множество всех изображений объекта, полученных при всевозможных условиях освещения, являющееся в данном случае \emph{формой изображения} объекта

Действительно, пусть изображение $f\in\R$ рассматриваемого объекта получено при условиях освещения $L$. Изменим в $\alpha\geqs0$ раз яркости всех источников света, формирующих освещение $L$, и обозначим новые условия освещения $\alpha L$. При изменении освещения $L$ на $\alpha L$ изменяется в $\alpha$ раз интенсивность падающего на объекта света, следовательно, в $\alpha$ раз изменяется и интенсивность отражённого света, а значит и световой поток, падающий на каждый детектор камеры. Поэтому, в силу условия, налагаемого на камеру, изображение объекта, полученное при освещении $\alpha L$, есть $\alpha f$. Таким образом, если $f$~--- изображение объекта, то $\alpha f$~--- изображение того же объекта, полученное, возможно, при других условиях освещения, $\alpha\geqs0$.
\newlength{\bearheight}%
\setlength{\bearheight}{0.18\textwidth}
\begin{figure}[htb!]
\centering
\begin{minipage}[b]{0.25\bearheight}
\scalebox{2}{$\alpha$}
\vspace*{0.40\bearheight}
\end{minipage}
\begin{minipage}[b]{0.3\bearheight}\centerline{\Large\bf \hspace{-5mm}$\times$}\vspace*{0.41\bearheight}\end{minipage}%
\includegraphics[height=\bearheight]{pictures/bear_alpha.png}
\begin{minipage}[b]{0.3\bearheight}\centerline{\Large\bf \hspace{-1mm}=}\vspace*{0.41\bearheight}\end{minipage}%
\includegraphics[height=\bearheight]{pictures/bear_alpha_2.png}
\begin{minipage}[b]{0.25\bearheight}
\scalebox{2}{$\;,\;\alpha \geq 0$}
\vspace*{0.10\bearheight}
\end{minipage}
\caption{$f - $изображение объекта. При умножении на константу $\alpha$ получается изображение того же объекта с измененной яркостью пикселей в $\alpha$ раз.}
\end{figure}

Рассмотрим теперь изображения $f_1\in\R$ и $f_2\in\R$ рассматриваемого объекта, полученные при условиях освещения $L_1$ и $L_2$ (наличие какой-либо связи между освещениями $L_1$ и $L_2$ в данном случае не предполагается). Пусть $L_{1+2}$ есть освещение, представляющее собой суперпозицию освещений $L_1$ и $L_2$, то есть $L_{1+2}$ получено одновременным включением источников света, формирующих освещения $L_1$ и $L_2$. В силу линейности и некогерентности источников света, световой поток, падающий на каждый детектор камеры при освещении $L_{1+2}$ есть сумма световых потоков, формируемых освещениями $L_1$ и $L_2$ в отдельности. Поэтому изображение объекта, полученное при освещении $L_{1+2}$, есть $f_1 + f_2$. Таким образом, если $f_1$ и $f_2$~--- изображения одного объекта, то $f_1 + f_2$~--- изображение того же объекта.
\begin{figure}[htb!]
\centering
\label{bear_real}%
\includegraphics[height=\bearheight]{pictures/bear_1-2}%
\begin{minipage}[b]{0.25\bearheight}\centerline{\Large\bf +}\vspace*{0.39\bearheight}\end{minipage}%
\includegraphics[height=\bearheight]{pictures/bear_2-2}%
\begin{minipage}[b]{0.3\bearheight}\centerline{\Large\bf =}\vspace*{0.41\bearheight}\end{minipage}%
\includegraphics[height=\bearheight]{pictures/bear_sum-2}%
\caption{Сумма двух изображений одного объекта, снятых при одинаковом ракурсе, но различном освещении, является изображением того же объекта с яркостью, равной сумме яркостей двух исходных изображений.}
\end{figure}

Суммируя вышесказанное, получаем, что форма изображения объекта, снимаемого при неизменном ракурсе $\vphi$, есть множество $V_\vphi\subset\R$, удовлетворяющее условиям
\begin{gather*}
    f\in V_\vphi \;\Rightarrow\; \alpha f \in V_\vphi, \ \alpha\geqs0,\\
    f_1,\, f_2\in V_\vphi \;\Rightarrow\; f_1 + f_2\in V_\vphi.
\end{gather*}
То есть форма изображения $V_\vphi$ является выпуклым конусом.

\subsection{Форма изображения объекта, регистрируемого при изменяющемся ракурсе}
\mbox{}\vspace{-\baselineskip}

Пусть теперь изменяться могут не только условия освещения объекта, но и ракурс съёмки $\vphi$. Множество, в пределах которого может принять своё значение векторный параметр $\vphi$, обозначим $\Phi$. \emph{Форма изображения} $V$ рассматриваемого объекта в этом случае может быть выражена следующим образом:
$$
    V = \bigcup_{\vphi\in\Phi} V_\vphi,
$$
где $V_\vphi$~--- выпуклый конус, представляющий собой множество изображений объекта, зарегистрированных при фиксированном ракурсе $\vphi\in\Phi$ и всевозможных условиях освещения. Такое множество $V$ является конусом, но, вообще говоря, не выпуклым (закрашенная серым область на рис. 3).

\begin{figure}[htb!]
\center\includegraphics[height=0.4\textwidth]{pictures/cone-non_convex.png}
\caption{Наглядное представление формы изображения объекта, снимаемого при всевозможных условиях освещения и ракурсах. Точками выделены выпуклые конусы $V_{\vphi_1}$ и $V_{\vphi_2}$, каждый из которых есть множество изображений объекта, зарегистрированных при одном ракурсе, но разных условиях освещения. Форма изображения объекта (выделена серым) есть объединение выпуклых конусов, соответствующих всевозможным ракурсам съёмки; она представляет собой конус, но, вообще говоря, не выпуклый. Её дополнение до выпуклого конуса (выделено штриховкой) состоит из изображений, которые не могут быть зарегистрированы в реальности}
\label{non_convex_cone}
\end{figure}


Обозначим $\ol V$ выпуклую оболочку формы $V$, которая в силу теоремы Каратеодори и того, что размерность $\dim\R = n$ пространства изображений $\R$ конечна, может быть представлена в виде
\begin{equation}
    \ol V = \left\{\sum_{i=1}^{n+1}\alpha_i f_i \;\middle|\; \alpha_i \geqs 0,\ f_i\in V,\ i=1,\, \ldots,\, n+1,\ \sum_{i=1}^{n+1}\alpha_i = 1  \right\},
\end{equation}
и в рассматриваемом случае является выпуклым конусом. Обозначим $\delta V$ дополнение формы $V$ до её выпуклой оболочки $\ol V$, т.\,е. $\delta V = \ol V \setminus V$. На рис. \ref{non_convex_cone} множество $\delta V$ обозначено штриховкой. Всякое изображение $f$ из $\delta V$ может быть представлено в виде выпуклой комбинации некоторых изображений рассматриваемого объекта, среди которых найдутся изображения, зарегистрированные при разных ракурсах и входящие в комбинацию с положительными коэффициентами (иначе $f\in V_\vphi\subset V$ при некотором $\vphi\in\Phi$). Подобные изображения выглядят как наложения принципиально разных изображений, в связи с чем далее будем считать, что они не являются изображениями каких-либо реальных объектов и не могут быть зарегистрированы.

Суммируя вышесказанное, получаем, что форма изображения объекта, снимаемого при произвольном ракурсе, есть конус $V$, выпуклая оболочка которого $\ol V$ не содержит никаких реальных изображений, кроме изображений рассматриваемого объекта.


\section{Искусственные нейронные сети}
Искусственная нейронная сеть ~--- математическая модель \cite{DudaHurt}, \cite{Bishop}, а также ее аппаратное воплощение, построеннаые по принципу организации и функционирования биологических нейронных сетей. Это понятие возникло при изучении процессов, протекающих в мозге, и при попытке смоделировать эти процессы. Структурной единицей является искусственный нейрон.

\emph{Искусственный нейрон} ~--- нелинейная функция $n$ аргументов $\varphi: \RR^n\to\RR$ вида $\varphi(x_1,\, \ldots,\, x_n) = f\left( \sum\limits_{i=1}^n w_ix_i \right)$, где $w_i$~--- весовые коэффициенты,\;\;\; $f$~--- функция активации. (рис. 4)
\begin{figure}[htb!]
\centering
\begin{tikzpicture}
  \draw[->] (6, -1) -- (7, -1) node[draw=none,fill=none,font=\scriptsize,midway,above] {$w_2$};
  \draw  (7, -1) -- (8, -1);
  \draw (8, -1) -- (6, -1) node[draw=none, fill=none, font=\scriptsize, above] {$x_2$};
  \draw (8, -1) -- (6, 1) node[draw=none, fill=none,font=\scriptsize, midway, above] {$w_1$};
  \draw (8, -1) -- (6, 1) node[draw=none, fill=none, font=\scriptsize, above] {$x_1$};
  \draw[->] (6, 1) -- (7, 0);
  \draw (7, 0) -- (8, -1);
  draw (8, -1) -- (6, -3) ;
  \draw (6, -3 )-- (8, -1) node[draw=none, fill=none, font=\scriptsize, midway, above] {$w_n$};
  \draw (8, -1 ) -- (6, -3) node[draw=none,fill=none,font=\scriptsize,above] {$x_n$};
  \draw[->] (6, -3) -- (7, -2);
  \draw[->] (8, -1)-- (9.2, -1) node[draw=none, fill=none, font=\scriptsize, below] {};
  \draw (9.2, -1) -- (9.5, -1);
 % \path[fill=red!40!yellow!40,draw=black] (8, -1) circle (8mm) ; 
 	\node[circle, align = center, text width=10mm, fill=blue!30] (all) at (8, -1) {   $f(\cdot)$};
\end{tikzpicture}



\caption{Модель нейрона: функция, принимающая на вход взвешенные $x_1, x_2, x_3$ (то есть $f(w_1 x_1 + w_2 x_2 + w_3 x_3)$), где $f$ ~-- нелинейная функция своих аргументов.}
\label{neuron}
\end{figure}
\hfill

\begin{figure}[htb!]
\centering
\def\layersep{3cm}
\begin{tikzpicture}[shorten >=1pt,->,draw=black!50, node distance=\layersep]
    \tikzstyle{every pin edge}=[<-,shorten <=1pt]
    \tikzstyle{neuron}=[circle,fill=black!25,minimum size=17pt,inner sep=0pt]
    \tikzstyle{input neuron}=[neuron, fill=green!50];
    \tikzstyle{output neuron}=[neuron, fill=red!50];
    \tikzstyle{hidden neuron}=[neuron, fill=blue!50];
    \tikzstyle{annot} = [text width=4em, text centered]

    % Draw the input layer nodes
    \foreach \name / \y in {1,...,4}
    % This is the same as writing \foreach \name / \y in {1/1,2/2,3/3,4/4}
        \node[input neuron, pin=left:Вход \#\y] (I-\name) at (0,-\y) {};

    % Draw the hidden layer nodes
    \foreach \name / \y in {1,...,5}
        \path[yshift=0.5cm]
            node[hidden neuron] (H-\name) at (\layersep,-\y cm) {};

    % Draw the output layer node
    \node[output neuron,pin={[pin edge={->}]right:Выход}, right of=H-3] (O) {};

    % Connect every node in the input layer with every node in the
    % hidden layer.
    \foreach \source in {1,...,4}
        \foreach \dest in {1,...,5}
            \path (I-\source) edge (H-\dest);

    % Connect every node in the hidden layer with the output layer
    \foreach \source in {1,...,5}
        \path (H-\source) edge (O);

    % Annotate the layers
    \node[annot,above of=H-1, node distance=1cm] (hl) {Скрытый слой};
    \node[annot,left of=hl] {Входной слой};
    \node[annot,right of=hl] {Выходной слой};
\end{tikzpicture}
\caption{Простая трехслойная нейронная сеть, имеющая 4 нейрона на входном слое, 5 нейронов на скрытом и еще один ~-- на выходном}
\label{simple_nn}
\end{figure}
Соединяя нейроны между собой и распологая их на разных слоях, можно сконструировать нейронную сеть (рис. \ref{simple_nn})

\newpage 
Современные нейросети могут содержать миллионы нейронов, расположенных на сотнях слоев и связанных друг с другом. Область применения искусственных нейронных сетей охватывает практически все сферы деятельности: прогнозирование курсов валют и цен на сырье, постановка диагнозов больным, кодирование и декодирование информации, оптимизация режимов производственного процесса, предсказание результатов выборов, анализ сейсмических данных, анализ сигналов (видео-, аудио- и проч.) и многие другие. 

Часто, сконструированная из некоторых соображений сеть оказывается настолько громоздкой, что нет никакой возможности строго математически объяснить, как и, главное, почему она решает поставленную задачу. Тем не менее за всей сложной структурой таких сетей стоит определенная математическая формула, которая входному элементу сети $X$ ставит в соответствие некоторое значение $Y$ ~--- выход сети. Именно этот факт и взят за основу данной работы.


\section{Постановка задачи}
Пусть даны изображения $f_1, f_2, \ldots, f_N$, а также изображение $g$. Требуется найти проекцию изображения $g$ на конус $V_F$~--- выпуклую оболочку  $f_1,\, \ldots,\, f_N$, т.\,е. 
$$V_F = \left\{ \sum\limits_{i=1}^N z_i f_i \;\middle|\; z_i \geqslant 0,\ i=1,\, \ldots,\, N \right\}$$.

Обозначим $F$ матрицу, элементы которой~--- координаты $f_1, f_2, \ldots, f_N$ в некотором ортонормированном базисе пространства изображений $\R$

Задача проецирования $g$ на конус $V_F$ -- это задача квадратичного программирования:

\begin{equation*}
	\begin{cases}
	\frac{1}{2}\norm{g - \sum\limits_{i=1}^N z_i f_i}^2 \sim 	\min\limits_{z_i} \\
	z_i \geq 0, \;\; i = 1, ..., N
	\end{cases}
\end{equation*}

Заметим, что $\norm{g - \sum\limits_{i=1}^N z_i f_i}^2 = \norm{g}^2 - 2 (g, F z) + \norm{F z}^2 \sim \min\limits_{z_i}\;$. Поскольку первое слагаемое $\norm{g}^2$ не зависит от $z$, а $\norm{Fz}^2 = (Fz,Fz) = (F^T F z, z)$, то эта задача минимизации эквивалента задаче $(F^T F z, z) - 2 (g, F z) \sim \min\limits_{z_i}$.

Тогда: $\frac{1}{2}\norm{g - \sum\limits_{i=1}^N z_i f_i}^2 = \frac{1}{2} (F^T F z, z) - (F^T g, z)$. Сделаем замену: $M = F^T F$ и получим более привычный вид задачи квадратичного программирования, которую и будем решать 
\begin{equation}\label{eq:qua_pro}
\frac{1}{2} (M z, z) - (F^T g, z) \sim \min\limits_{z \geq 0}.
\end{equation}



\section{Метод проективных градиентов в задаче проецирования на многогранный конус}
\subsection{Итерационный алгоритм метода}
Для поставленной задачи квадратичного программирования \eqref{eq:qua_pro} был выбран итерационный алгоритм \cite{LCP}, имеющий некоторые особенности: он предполагает положительную определьнность матрицы $M$, а также ее симметричность. Обозначим $q = - F^T g$.
$$
        z^{r+1} = \sigma [z^r - \omega E^r (M z^r + q + K^r(z^{r+1} - z^r))]^+ +(1 - \sigma ) z^r
    $$
    Имея некий выбор при подборе параметров, положим $K^r = 0, \; \sigma = 1, \; E^r = I$, получим: 
    \begin{equation}\label{eq:proj_grad}
    z^{r+1} =[(z^r - \omega (M z^r + q)]^+ = [(I - \omega M) z^r - \omega q]^+ 
    \end{equation}
        
Выбор $\omega$ ограничивается условием:
    $$[2 (\omega I)^{-1} - M] \mbox{ ~--- положительно определенная матрица}$$
    Или, что то же самое, $$\min\lambda(2 (\omega I)^{-1}) > \max\lambda (M) \Leftrightarrow 0 < \omega < \frac{2}{\max\lambda(M)},$$
    где $\lambda(2 (\omega I)^{-1})$ и $\lambda(M)$~--- спектры матриц $2 (\omega I)^{-1}$ и $M$.
    
\subsection{Нейросетевая реализация метода}
Формула \eqref{eq:proj_grad} такой структуры реализуема в виде рекуррентной нейронной сети, обладающей некоторыми важными отличительными чертами, делающими подобные сети чрезвычайно интересными для исследований (но при этом и заметно усложняющими алгоритмы и методы работы с ними). Такие сети имеют одну или несколько обратных связей между различными слоями нейронов, что позволяет запоминать и воспроизводить последовательности реакций на входной импульс. С точки зрения программирования такие сети являются аналогом циклов.

\begin{figure}[htb!]
\centering
\begin{tikzpicture}[->,>=stealth']
    \node (Input)  [ellipse, text width=30mm, align=center, fill=blue!30] {Input \;\;$z$};
    \node (n) [circle, above=of Input.north, yshift=-0.2cm, text width=3mm, align=center, fill = blue!30] {+};
    \node (Hidden) [rectangle, above=of n.north, align=center, yshift = -0.4cm, text width = 45mm, fill = blue!30] {$z^{r+1}\;$  Recurrent layer $f(x) = max\{0, x\}$ };
    \node (e) [right=of Hidden.south, yshift=-3.3em] {$(I - \omega F^T F)$};
    \node (Output) [ellipse, above=of Hidden, yshift=0.2cm, align = center, text width = 30mm, fill = blue!30] {Output \;\;$z_*(g)$};
   
    \draw [->]  (Input)  edge node[anchor = east] {$\omega F^T$}  (n);
    \draw [->] (n) edge (Hidden);
    \draw [->]  (Hidden)  edge node[anchor = west] {} (Output);
    \path [line,rounded corners] ($(Hidden.north) + (+0.0, 0.1)$) |- ($(Hidden.south) + (3.5, 1.8)$) |- (n);
\end{tikzpicture} 
\caption{Структура реализующей итерационный метод проективных градиентов рекуррентной сети. При обучении такой сети со связанными между собой весовыми матрицами стандартные алгоритмы градиентного спуска не сходятся при попытке вмешательства в их работу, нацеленной на поддержание накладываемых методом связей. Поэтому возникает потребность в доработке данной нейронной сети.}
\label{first_rnn}
\end{figure}
Изначально была построена нейронная сеть со структурой, показанной на рис. \ref{first_rnn}. Корректность работы такой сети была протестирована на различных изображениях разных объектов. 

Согласно итерационному алгоритму, проекция будет равна $\Pi_{V_g} g = F z$, что полностью соответствует нашим представлениям о рассматриваемой модели изображения: например, если в матрицу $F$ записать изображения объекта $f_1,...,f_N$, а потом подать на ее вход изображение $f_i$, то получаемая на выходе матрица $z$ будет иметь следующий вид: $z = (0, ..., 0, 1, 0, ...0)$, где 1 стоит на $i$-позиции. Это означает, что если задать <<грани>> конуса $f_1,...,f_N$, который является формой изображения рассматриваемого объекта, а затем подать на вход сети изображение $f_i$, соответствующее одной из таких граней, то проекция $F z = f_i$ совпадает с $f_i$, что и является ожидаемым результатом: проекции всех изображений, принадлежащих конусу, должны совпадать с самими изображениями.

Заметим, что указанные на рис. \ref{first_rnn} весовые матрицы связаны между собой параметром $F^T$. Этот факт стал главной проблемой при попытке обучать такую сеть (то есть задав весовые матрицы произвольно, попытаться <<научить>> сеть выполнять поставленную задачу~-- проецирование на конус) в виду необходимости поддерживать эту связь. Была предпринята попытка после каждого этапа обучения вручную подправлять параметры сети, чтобы они удовлетворяли условию связи. Однако, стандартные методы обучения методом градиентного спуска не срабатывали, так как такое корректирование параметров может направлять нас в сторону, противоположное градиенту спуска ~--- алгоритм не сходился.

Таким образом, появилась необходимость переосмыслить структуру весовых матриц и накладываемые на них условия.
  
\subsection{Интерпретация возможной нейросетевой реализации алгоритма}

Пусть $M$ ~--- положительно определенная симметричная матрица. Тогда $$(M z, z) - 2 (F^T g, z) = (M^{1/2} z, M^{1/2} z) - 2 (M^{-1/2} F^T g, M^{1/2} z) = \norm{M^{1/2} z - M^{-1/2} F^T g}^2$$

Таким образом, мы приходим к эквивалентной задаче: 
\begin{equation}\label{eq:final_min_problem}
\norm{M^{1/2} z - M^{-1/2} F^T g}^2 + \norm{g}^2 - \norm{F (F^T F)^{-1} F^T g}^2 \sim \min\limits_{z \geq 0}
\end{equation}

Решение задачи (\ref{eq:final_min_problem}) обозначим $z_*(g)$. Оно может быть найдено методом проективных градиентов, а значит, и с использованием нейросетевой реализации этого метода.

Далее, обозначим 
$$V_{F, M} \myeq \{F (F^T F)^{-1} M z,\;z \geq 0\}$$

Легко показать, что  $V_{F, M}$ ~--- выпуклый конус:

1) если $g \in V_{F, M}$, \;то $\alpha g \in V_{F, M}, \; \forall \alpha \geq 0$  

2) если $g_1,\;g_2 \in V_{F, M}$, то $g_1 + g_2 \in V_{F, M}$

Введем еще одно обозначение: $$\Pi_{F,M} \; g \myeq F (F^T F)^{-1} M z_*(g)$$
Заметим, что $F^T \Pi_{F, M} g = F^T F (F^T F)^{-1} M z_*(g) = M z_*(g)$.

Следовательно, $\Pi_{F, M}^2 g = \Pi_{F, M} g = F (F^T F)^{-1} M z_*(g) \;\;$, то есть при повторном применении этот оператор дает тот же результат, что и при одинарном. Это означает, что $\Pi_{F, M}$~-- идемпотентный оператор, построенный на основе нейросети, представленной на рис. \ref{first_rnn}.


\subsection{Полученная структура нейросети}
В итоге приходим к такой структуре нейросети:
\vspace{1cm}
\begin{figure}[h!]

{\center
\begin{tikzpicture}[->,>=stealth']
    \node (Input)  [ellipse, text width=20mm, align=center, fill=blue!30] {Вход \;\;g};
    \node (n) [circle, above=of Input.north, yshift=-0.2cm, text width=3mm, align=center, fill = blue!30] {+};
    \node (Hidden) [rectangle, above=of n.north, align=center, yshift = -0.4cm, text width = 45mm, fill = blue!30] {Рекуррентный слой $z^{r+1}\;$  $f(x) = \max\{0, x\} - \text{ReLU}$ };
    \node (e) [right=of Hidden.south, yshift=-3.2em] {$(I - M)$};
    \node (Output) [ellipse, above=of Hidden, yshift=0.2cm, align = center, text width = 20mm, fill = blue!30] {Выход \;\;$z_*$};
    
    
    \node (comment1) [rectangle, text width = 0.5\linewidth, left = of Hidden, xshift = +2.5em, yshift = -2em]{%
    \begin{flushleft}
    \vspace{-1cm}\hspace{-0.8cm}$\norm{M^{1/2} z - M^{-1/2}F^Tg}^2 \sim \min\limits_{z \geq 0}$\\
    \vspace{0.3cm}\hspace{-0.8cm}$\Pi_{F, M} g = F (F^T F)^{-1} M z_*(g)$\\
	\vspace{0.3cm}\hspace{-0.8cm}$V_{F, M} = \{F (F^T F)^{-1} M z,\;z \geq 0\}$\\    
	\vspace{0.3cm}\hspace{-0.8cm}$\Pi_{F, M} g = arg \min\limits_{f \in V_{F, M}} \norm{M^{-1/2} F^T (g - f)}^2$\\
	\vspace{0.7cm}\hspace{-0.8cm} Ограничения на $M:$\\
	$M, (2 I - M) - $ неотр. определены \\
	$M = M^T$
	\end{flushleft}}; 
   
    \draw [->]  (Input)  edge node[anchor = east] {$F^T$}  (n);
    \draw [->] (n) edge (Hidden);
    \draw [->]  (Hidden)  edge node[anchor = west] {} (Output);
    \path [line,rounded corners] ($(Hidden.north) + (-0.0, 0.1)$) |- ($(Hidden.south) + (3.5, 2.4)$) |- (n); 
\end{tikzpicture}}
\caption{Схема полученной структуры рекуррентной нейросети.}
\label{schema:final_net}
\end{figure}\\

Важно отметить, что
$$\norm{M^{1/2} z - M^{-1/2} F^T g}^2 + \norm{g}^2 - \norm{F (F^T F)^{-1} F^T g}^2 \; = 0 \;\;\Leftrightarrow \;\; g \in V_{F, M}$$


\section{Обучение нейросети}
\subsection{Задание функции потерь}

Для обучения сети используется тестовая выборка объема $L$~-- набор $L$ изображений, которые подаются на вход нейронной сети. С помощью получаемых на ее выходе $z$ считается функция потерь loss следующим образом

\begin{equation}\label{eq:loss}
loss = \sum\limits_{i=1}^L s_i\; \left(\;\norm{ M^{-1/2} (F^T g^{(i)} - M z_*(g^{(i)})) }^2 + \norm{g}^2 - \norm{(F^T F)^{-1/2} F^T g^{(i)}}^2\right)
\end{equation}


\hspace{1.5cm}
$$
	s_i = 
	\begin{cases}
		\;\;\;1, \;\; g^{(i)} \in V_{F, M}\\
		-1, \;\; g^{(i)} \notin V_{F, M}
	\end{cases}, \;\;i = 1, ..., L
$$

Реализация нейронной сети была осуществлена на языке программирования Python с помощью Theano~-- библиотеки Python, которая позволяет определять, оптимизировать и рассчитывать математические выражения эффективно используя символьные вычисления. Среди возможностей этой библиотеки стоит особо отметить быструю и стабильную оптимизацию, а также эффективное дифференцирование сложных функций:
$$\frac{\partial f}{\partial x} = \sum\limits_i \frac{\partial f}{\partial y_i} \frac{\partial y_i}{\partial x}.$$ Тем не менее, заданную функцию потерь loss по формуле (\ref{eq:loss}) theano дифференцировать не умеет, поэтому отдельной задачей, которую удалось решить, была реализация правильного градиента для такой функции. Поскольку в Theano выражения задаются в символьном виде, а встроенный оптимизатор упрощает получившиеся градиенты и выражения, численное значение которых требуется получить, мы получаем некоторый прирост производительности вычислений. 

\subsection{Корректировка матриц весовых коэффициентов}
При этом само обучение проводилось модифицированным методом обратного распространения ошибки: помимо стандартной процедуры следовало учесть особенности и ограничения, накладываемые используемым методом проекции градиента ~-- ограничения на матрицу $M$:
\begin{itemize}
\item $M, (2 I - M) - $ неотр. определены
\item $M = M^T$
\end{itemize}

Для этого производится корректировка параметров сети: после каждой итерации обучения вычисляется приращение $\widetilde{\delta M}$ для матрицы весовых коэффициентов $M$ в целях минимизации функции потерь (\ref{eq:loss}). Далее это приращение симметризуется, а собственные значения получившейся матрицы удерживаются в определенном диапазоне, чтобы удовлетворить накладываемым ограничениям:\\

$\widetilde{\delta M} = \frac{1}{2} (\delta M + \delta M^T)$

$S, V = eig(M + \widetilde{\delta M})$

$\widetilde{S_i} = \max(-0.95, \;\min(S_i,\; 0.95))$

Окончательное значение для $M$: $M = V \widetilde{S} \;V^T$


\subsection{Обучающая и тестовая выборки}
Обучение и тестирование нейросети производилось на изображениях двух объектов: медвежонка и девушки, снятых при различных условиях регистрации различных условиях освещения и ракурса съемки)
\begin{figure}[htb!]
\centering
\includegraphics[height=0.15\textwidth]{pictures/1}\hspace{0.05mm}
\includegraphics[height=0.15\textwidth]{pictures/2}\hspace{0.05mm}
\includegraphics[height=0.15\textwidth]{pictures/3}\hspace{0.05mm}
\includegraphics[height=0.15\textwidth]{pictures/4}

\vspace{1mm}
\includegraphics[height=0.15\textwidth]{pictures/18}\hspace{0.05mm}
\includegraphics[height=0.15\textwidth]{pictures/12}\hspace{0.05mm}
\includegraphics[height=0.15\textwidth]{pictures/16}\hspace{0.05mm}
\includegraphics[height=0.15\textwidth]{pictures/14}
\caption{Примеры изображений двух объектов, используемых для обучения и тестирования нейросети: девушки и медвежонка, снятые при различных условиях освещения и ракурсах}
\end{figure}

\section{Результаты}
В ходе выполнения данной работы были получены следующие результаты: 

\begin{itemize}
\item Реализованы два типа нейронных сетей:
	\begin{itemize}
	\item[--] сеть со связанными матрицами весовых коэффициентов, реализующая проецирование на форму изображения как многогранный конус
	\item[--] сеть с независимыми матрицами весовых коэффициентов, позволяющая идентифицировать объекты по их изображениям
	\end{itemize}
\item Реализованы функция потерь и ее градиент для указанных сетей, что позволило использовать для их обучения модифицированные методы обратного распространения ошибки
\item Проведен вычислительный эксперимент с изображениями реальных объектов
\end{itemize}

В вычислительном эксперименте были получены следующие численные знаничения точности идентификации (здесь \emph{число итераций}~-- количество <<проходов>> по рекуррентному слою (см. рис. \ref{schema:final_net}) до того, как получаемый на выходе результат будет считаться окончательным, или, что то же самое, число итераций в формуле метода проективных градиентов (\ref{eq:proj_grad}))\\

\begin{minipage}{.50\textwidth}

\textbf{\textit{\Large Без обучения}}
\begin{itemize}
\item \textbf{\textit{Число итераций: 100}}
\begin{itemize}
\item Изображения одного ракурса:\\ \textbf{90-95\%}
\item Изображения разных ракурсов:\\ \textbf{80-85\%}
\end{itemize}
\item \textbf{\textit{Число итераций: 500}}
\begin{itemize}
\item Изображения одного ракурса:\\ \textbf{95-100\%}
\item Изображения разных ракурсов:\\ \textbf{85-95\%}
\end{itemize}
\end{itemize}
\end{minipage}
\begin{minipage}{.50\textwidth}
\vspace{-2.4cm}\textbf{\textit{\Large С обучением}}
\begin{itemize}
\item \textbf{\textit{Число итераций: 100}}
\begin{itemize}
\item Изображения одного ракурса:\\ \textbf{80-85\%}
\item Изображения разных ракурсов:\\ \textbf{70-80\%}
\end{itemize}

\item \textbf{\textit{Число итераций: 500}}

\centering
\textbf{\textit{\Large ???}}
\end{itemize}
\end{minipage}\\\\

В силу ограниченной мощности вычислительной машины, на которой производилось тестирование, нет возможности провести обучение сети при числе итераций 500. Тем не менее это представляет большой исследовательский интерес.


\newpage
\section{Литература}

\begin{thebibliography}{10}

\bibitem{PitievChulichkov_morph_methods} Ю. П. Пытьев, А. И. Чуличков <<Методы морфологического анализа изображений>>, ФИЗМАТЛИТ, 2010

\bibitem{Pitiev_Zadachi_Morph_Analiza} Ю.П. Пытьев <<Задачи морфологического анализа изображений>>, 1984

\bibitem{TuGonsales} Дж. Ту, Р. Гонсалес <<Принципы распознавания образов>>, Издательство <<Мир>>,\; Москва 1978

\bibitem{DudaHurt} Дуда Р., Харт П. <<Распознавание образов и анализ сцен>>, Издательство <<МИР>>, Москва 1976

\bibitem{ZubukFedotov} Зубюк А.В., Федотов А.Б. <<Алгоритм идентификации объектов по изображениям, основанный на разделении гиперплоскостью и нечувствительный к изменению условий освещения и ракурса>>. Вестник Московского университета, 2015

\bibitem{LCP} Katta G. Murty, <<Linear Complementarity, Linear and Nonlinear Programming>>, Internet Edition

\bibitem{Bishop} Christopher M. Bishop <<Pattern Recognition and Machine Learning>>, Springer, 2006

\end{thebibliography}

\end{document}
